\begin{thema}{Γ}
Στο σχήμα δίνεται η γραφική παράσταση $C_f$ μιας δύο φορές παραγωγίσιμης συνάρτησης $f$ ορισμένης στο $[0, 6]$. Είναι γνωστό ότι $f(6) = -2$.

\begin{center}
\begin{tikzpicture}[scale=0.7]
  \draw[help lines, gray!30] (-1,-3) grid (7,4);
  % curve
  \draw[very thick, maincolor, smooth, tension=0.55]
    plot coordinates {(0,0) (0.5,1) (1,2.2) (1.5,2.8) (2,3) (2.5,2.8) (3,2) (3.5,1) (4,0) (5,-1.5) (6,-2)};
  % shaded positive area
  \fill[maincolor!12, opacity=0.5, smooth, tension=0.55]
    (0,0) -- plot coordinates {(0,0) (0.5,1) (1,2.2) (1.5,2.8) (2,3) (2.5,2.8) (3,2) (3.5,1) (4,0)} -- cycle;
  % shaded negative area
  \fill[maincolor!10, opacity=0.5, smooth, tension=0.55]
    (4,0) -- plot coordinates {(4,0) (5,-1.5) (6,-2)} -- (6,0) -- cycle;
  \filldraw[maincolor] (0,0) circle (2pt);
  \filldraw[maincolor] (2,3) circle (2pt) node[above,font=\footnotesize] {$(2,3)$};
  \filldraw[maincolor] (4,0) circle (2pt);
  \filldraw[maincolor] (6,-2) circle (2pt) node[below right,font=\footnotesize] {$(6,-2)$};
  \node at (4.5,2) {$C_f$};
  \draw[-latex] (-1.3,0) -- (7.3,0) node[right] {$x$};
  \draw[-latex] (0,-3.3) -- (0,4.3) node[above] {$y$};
\foreach \x in {1,2,3,4,5,6}
    \draw (\x,0.1) -- (\x,-0.1) node[below,font=\footnotesize] {$\x$};
  \foreach \y in {-2,-1,1,2,3}
    \draw (0.1,\y) -- (-0.1,\y) node[left,font=\footnotesize] {$\y$};
\end{tikzpicture}
\end{center}

\begin{erwthma}
  \item Να μελετήσετε την $f$ ως προς τη μονοτονία, να βρείτε τα τοπικά της ακρότατα και να προσδιορίσετε το σύνολο τιμών της. Στη συνέχεια, να βρείτε το πλήθος των ριζών της εξίσωσης $f(x) = 1$.
  \item Να ελέγξετε αν ικανοποιούνται οι προϋποθέσεις του Θεωρήματος \eng{Rolle} για την $f$ στο διάστημα $[0, 4]$ και να διατυπώσετε το γεωμετρικό συμπέρασμα του θεωρήματος.
  \item Να αποδείξετε γεωμετρικά ότι $\dint_0^4 f(x)\,\d x > 6$.
  \item Αν επιπλέον γνωρίζετε ότι η εφαπτομένη της $C_f$ στο σημείο $O(0,0)$ έχει εξίσωση $y=2x$, να υπολογίσετε το όριο:
  \[ \lim\limits_{x \to 0} \frac{f(x) - \hm x}{x^2 + x} \]
\end{erwthma}
\end{thema}
